\section{Altre forme di limiti e colimiti}
Lo scopo di questa sezione è generalizzare le costruzioni date precedentemente a limiti e colimiti di forma diversa da prodotti, somme, equalizzatori e coequalizzatori.

Iniziamo coi prodotti di taglia arbitraria (piccola):

esempi grandi: in Set, Top, k-Vect, gruppi, anelli etc., Cat, e molte altre categorie del cap 1

esempi piccoli: un (pre)ordine, (ordini, insiemi di parti etc...)

non esempi: categorie senza prodotti, prodotti non categoriali (tensore --che comunque ha una proprietà universale)

unicità

prodotto di famiglie

oggetti terminali.

equalizzatori e coequalizzatori (coppie nucleo e il quoziente di una relazione di equivalenza interna, fattorizzazione di un morfismo)

prodotti fibrati (pullback) e somme amalgamate (pushout) (sottoalgebre definite da equazioni, amalgame)

colimiti filtrati

categorie puntate, nuclei e conuclei

esempi e non esempi

perché si chiamano limiti?
\begin{remark}
	\Todo{Limiti in analisi come limiti categoriali}
\end{remark}
\subsection{Limiti in algebra, topologia e logica}
\begin{example}
	Il \emph{nucleo} di una applicazione lineare \(f \colon V \to W\) tra spazi vettoriali (o tra gruppi abeliani, o \(R\)-moduli,\dots), già introdotto in \ref{ker_e_coker}, è un esempio di limite: si tratta dell'equalizzatore della coppia
	% https://q.uiver.app/#q=WzAsMixbMCwwLCJWIl0sWzEsMCwiVyJdLFswLDEsIjAiLDAseyJvZmZzZXQiOi0yfV0sWzAsMSwiZiIsMix7Im9mZnNldCI6MX1dXQ==
	\[\begin{tikzcd}[cramped]
			V & W
			\arrow["0", shift left=2, from=1-1, to=1-2]
			\arrow["f"', shift right, from=1-1, to=1-2]
		\end{tikzcd}\]
	In effetti una nozione di `nucleo' in questo senso si può dare in ogni categoria che abbia un \emph{oggetto zero}, nel senso di un oggetto che sia allo stesso tempo iniziale e terminale (questo è il caso di ciascuna delle categorie appena menzionate, ma anche della categoria degli insiemi puntati di \ref{}). In tal caso, si ricordi che la `mappa zero' consiste della composizione \(0_{VW} \colon V \to 0 \to W\) e in un certo senso evidente \(\ker f\) è in quel caso il massimo sottoggetto \(i_{\ker f} \colon \ker f \mono V\) tale che \(f\circ i_{\ker f} = 0_{VW}\). Nel caso degli insiemi puntati, dove \(V=(A,a_0)\), si tratta del sottoinsieme \(K\subseteq A\) determinato come \(K=\{a\in A\mid fa = b_0\}\), se \(W=(B,b_0)\) è l'insieme puntato codominio.
\end{example}
\begin{example}
	Il funtore dei sottoggetti introdotto in \ref{} fa uso della nozione di prodotto fibrato in \(\ctSet\): infatti, data una funzione \(u \colon A\to B\) tra insiemi, la funzione \(\Sub(u) \colon \Sub(B) \to \Sub(A)\) si determina dal prodotto fibrato
	\[\begin{tikzcd}
			u^*E =\pull AEB \arrow[r] \arrow[d, hook] & E \arrow[d, hook] \\
			A \arrow[r, "u"'] & B
		\end{tikzcd}\]
	Infatti, la funzione \(\Sub(u)\) manda un sottoggetto \(E\mono B\) nel sottoggetto \(u^*E\mono A\).
\end{example}
\begin{example}
	\Todo{Punti fissi come limiti}
\end{example}
\begin{example}
	\Todo{La categoria degli elementi come limite}
\end{example}
\begin{example}
	\Todo{\((F/G)\) come limite}
\end{example}
\begin{example}
	\Todo{\(Nat(F,G)\) come limite}
\end{example}
\begin{example}
	\Todo{Spazi di lacci come limite}
\end{example}
\begin{example}
	\Todo{Costruzione dei limiti in \(\ctCat\)}
\end{example}
\begin{example}
	\Todo{Estensione destra di F lungo G}
\end{example}
\begin{examples}
	Esempi di estensioni destre
\end{examples}
\subsection{Colimiti in algebra, topologia e logica}
\begin{example}
	\Todo{Prodotto semidiretto di gruppi e monoidi}
\end{example}
\begin{example}
	\Todo{}
\end{example}
\begin{example}
	\Todo{Conuclei e quozienti}
\end{example}
\begin{example}
	\Todo{Spazi di orbite come colimiti}
\end{example}
\begin{example}
	\Todo{\(\pi_0\) come colimite.}
\end{example}
\begin{example}
	\Todo{La categoria degli elementi come colimite}
\end{example}
\begin{example}
	\Todo{\(\ctC\star\ctD\) come colimite}
\end{example}
\begin{example}
	\Todo{Sospensione e Sospensione ridotta come colimite}
\end{example}
\begin{example}
	\Todo{Costruzione dei colimiti in \(\ctCat\)}
\end{example}
\begin{example}
	\Todo{Estensione sinistra di F lungo G}
\end{example}
\begin{examples}
	Esempi di estensioni sinistre
\end{examples}
\begin{theorem}
	Teorema di Knaster-Tarski (l'esistenza di un punto fisso è determinata da una `prop univ')
\end{theorem}
